\newif\ifuseseminar
\useseminartrue
\input{../../latex_common/header.tex.frag}

\title{Mecânica Quântica II -- Partículas Idênticas e Transformações}

\begin{document}
\begin{enumerate}
	\item Considere um sistema de duas partículas {\bf idênticas} e responda às questões
	      abaixo:
	      \begin{enumerate}
		      \item Utilize o produto tensorial para definir os dois possíveis espaços
		            para os estados de duas partículas.
		      \item Os estados representando duas partículas idênticas podem ser
		            separáveis? Justifique sua resposta.
		      \item Discuta os casos onde estados de duas partículas idênticas são
		            aproximadamente separáveis.
		      \item Explique como podemos testar empiricamente se o estado de duas
		            partículas é simétrico ou anti-simétrico. Descreva como a simetria e a
		            anti-simetria estão relacionadas ao conceito de troca de partículas
		            idênticas e como as técnicas experimentais podem ser usadas para medir os
		            estados de duas partículas e verificar sua simetria.
	      \end{enumerate}
	\item Considere um sistema uni-dimensional com uma partícula e responda às questões
	      abaixo:
	      \begin{enumerate}
		      \item Defina o operador de translação $\hat{T}_\epsilon$ que desloca a
		            função de onda de uma partícula de uma distância $\epsilon$.
		            Isto é, $\hat{T}_\epsilon\ket{x} = \ket{x+\epsilon}$.
		      \item Mostre que o gerador do operador de translação é o operador momento
		            linear $\hat{p}$. Qual é o gerador se permitirmos translações da
		            forma $\hat{T}_{\epsilon,g}\ket{x} = e^{i\epsilon
					            g(x)/\hbar}\ket{x+\epsilon}$?
	      \end{enumerate}
	\item Em um sistema tridimensional com uma partícula, responda às questões abaixo:
	      \begin{enumerate}
		      \item Defina o operador de reflexão $\hat{R}_x$ que inverte a coordenada
		            $x$. Mostre que o operador de reflexão é unitário e auto-adjunto.
		            Encontre as auto-funções e auto-valores do operador de reflexão.
		      \item Mostre que o operador paridade $\hat{P}$ pode ser escrito como
		            uma rotação seguida de uma reflexão. Mostre que o operador paridade é
		            unitário e auto-adjunto. Encontre as auto-funções e auto-valores do operador
		            paridade.
	      \end{enumerate}
	\item Inversão temporal é uma transformação que inverte o tempo. Considere um
	      sistema com uma partícula e responda às questões abaixo:
	      \begin{enumerate}
		      \item Usando a equação de Schrödinger, mostre que a inversão temporal
		            requer que $\psi(x,t)\rightarrow\psi^*(x,-t)$.
		      \item Mostre que a inversão temporal transforma a Hamiltoniana na
		            representação de Schrödinger em seu complexo conjugado, i.e.,
		            $\hat{H}(t)\rightarrow\hat{H}^*(-t)$.
		            O que isso implica sobre a evolução temporal do sistema?
	      \end{enumerate}
\end{enumerate}


\end{document}
